% ============================================================
% Chapitre 5 : Intervalles de Confiance
% ============================================================
\chapter{Intervalles de Confiance}
\label{ch:ic}

\begin{center}
\textit{<<~Une estimation sans mesure d'incertitude est une estimation incompl\`ete.~>>}
\end{center}

% ------------------------------------------------------------
\section{Exemple motivant}
% ------------------------------------------------------------

Un sondage interroge $n=1000$ personnes : 520 sont favorables \`a une r\'eforme. On estime la proportion $p$ par $\phat=0{,}52$. Mais quelle est la \emph{pr\'ecision} de cette estimation ? Un intervalle de confiance r\'epond : <<~avec 95\% de confiance, $p\in[0{,}489;\,0{,}551]$~>>.

% ------------------------------------------------------------
\section{D\'efinition et interpr\'etation}
% ------------------------------------------------------------

\begin{definition}[Intervalle de confiance]
\index{intervalle de confiance}
Un \textbf{intervalle de confiance} (IC) de niveau $1-\alpha$ pour $\theta$ est un intervalle al\'eatoire $[L(X_1,\ldots,X_n),\;U(X_1,\ldots,X_n)]$ tel que :
\[
\Prob_\theta(L\leq\theta\leq U) \geq 1-\alpha \quad \forall\theta\in\Theta.
\]
\end{definition}

\begin{attention}
L'interpr\'etation correcte est \textbf{fr\'equentiste} : si on r\'ep\`ete l'exp\'erience un grand nombre de fois, $(1-\alpha)\times 100\%$ des IC construits contiendront la vraie valeur $\theta$.

\medskip
\textbf{Erreur courante :} <<~Il y a 95\% de chances que $\theta$ soit dans l'IC~>>. Non ! $\theta$ est fix\'e ; c'est l'intervalle qui est al\'eatoire. Apr\`es observation, l'IC contient $\theta$ ou non.
\end{attention}

\begin{definition}[Quantit\'e pivotale]
\index{quantite pivotale@quantit\'e pivotale}
Une \textbf{quantit\'e pivotale} est une fonction $Q(X_1,\ldots,X_n,\theta)$ dont la loi ne d\'epend pas de $\theta$. On inverse la relation $\Prob(a\leq Q\leq b)=1-\alpha$ pour obtenir l'IC.
\end{definition}

% ------------------------------------------------------------
\section{IC pour la moyenne}
% ------------------------------------------------------------

\subsection{Variance connue}

\begin{theoreme}
Soit $X_1,\ldots,X_n\iid\mathcal{N}(\mu,\sigma^2)$ avec $\sigma^2$ connu. Un IC de niveau $1-\alpha$ pour $\mu$ est :
\[
\left[\xbar - z_{\alpha/2}\frac{\sigma}{\sqrt{n}},\;\xbar + z_{\alpha/2}\frac{\sigma}{\sqrt{n}}\right],
\]
o\`u $z_{\alpha/2}$ est le quantile d'ordre $1-\alpha/2$ de $\mathcal{N}(0,1)$.
\end{theoreme}

\begin{proof}
La quantit\'e pivotale est $Z=\frac{\xbar-\mu}{\sigma/\sqrt{n}}\sim\mathcal{N}(0,1)$. On a $\Prob(-z_{\alpha/2}\leq Z\leq z_{\alpha/2})=1-\alpha$, et on isole $\mu$.
\end{proof}

\subsection{Variance inconnue}

\begin{theoreme}
\label{thm:ic-student}
Soit $X_1,\ldots,X_n\iid\mathcal{N}(\mu,\sigma^2)$ avec $\sigma^2$ inconnu. Un IC de niveau $1-\alpha$ pour $\mu$ est :
\[
\left[\xbar - t_{\alpha/2,n-1}\frac{S}{\sqrt{n}},\;\xbar + t_{\alpha/2,n-1}\frac{S}{\sqrt{n}}\right],
\]
o\`u $t_{\alpha/2,n-1}$ est le quantile d'ordre $1-\alpha/2$ de la loi $t(n-1)$.
\end{theoreme}

\begin{proof}
La quantit\'e pivotale est $T=\frac{\xbar-\mu}{S/\sqrt{n}}\sim t(n-1)$ (cf.\ th\'eor\`eme~\ref{thm:loi-s2}).
\end{proof}

\begin{remarque}
Pour $n$ grand ($n\geq 30$ en pratique), $t_{\alpha/2,n-1}\approx z_{\alpha/2}$ et les deux IC co\"incident. De plus, par le TCL, l'IC de Student est approximativement valide m\^eme sans normalit\'e pour $n$ grand.
\end{remarque}

% ------------------------------------------------------------
\section{IC pour la variance}
% ------------------------------------------------------------

\begin{theoreme}
Sous le mod\`ele normal, un IC de niveau $1-\alpha$ pour $\sigma^2$ est :
\[
\left[\frac{(n-1)S^2}{\chi^2_{\alpha/2,n-1}},\;\frac{(n-1)S^2}{\chi^2_{1-\alpha/2,n-1}}\right],
\]
o\`u $\chi^2_{\alpha/2,n-1}$ et $\chi^2_{1-\alpha/2,n-1}$ sont les quantiles de la loi $\chi^2(n-1)$.
\end{theoreme}

\begin{proof}
La quantit\'e pivotale est $V=\frac{(n-1)S^2}{\sigma^2}\sim\chi^2(n-1)$.
\end{proof}

% ------------------------------------------------------------
\section{IC pour une proportion}
% ------------------------------------------------------------

\begin{theoreme}
Soit $X_1,\ldots,X_n\iid\mathcal{B}(1,p)$. Pour $n$ grand, un IC approximatif de niveau $1-\alpha$ pour $p$ est :
\[
\left[\phat - z_{\alpha/2}\sqrt{\frac{\phat(1-\phat)}{n}},\;\phat + z_{\alpha/2}\sqrt{\frac{\phat(1-\phat)}{n}}\right],
\]
o\`u $\phat=\xbar$.
\end{theoreme}

\begin{proof}
Par le TCL, $\frac{\phat-p}{\sqrt{p(1-p)/n}}\xrightarrow{\mathcal{L}}\mathcal{N}(0,1)$. On remplace $p$ par $\phat$ dans l'\'ecart-type (m\'ethode de Wald).
\end{proof}

\begin{attention}
L'IC de Wald peut \^etre mauvais pour $p$ proche de 0 ou 1, ou pour $n$ petit. L'IC de Wilson (ou Agresti--Coull) est pr\'ef\'erable :
\[
\tilde{p} = \frac{\phat n + z^2/2}{n+z^2}, \quad \tilde{p}\pm z_{\alpha/2}\sqrt{\frac{\tilde{p}(1-\tilde{p})}{n+z^2}}.
\]
\end{attention}

% ------------------------------------------------------------
\section{IC pour la diff\'erence de deux moyennes}
% ------------------------------------------------------------

\begin{theoreme}
Soient $X_1,\ldots,X_{n_1}\iid\mathcal{N}(\mu_1,\sigma^2)$ et $Y_1,\ldots,Y_{n_2}\iid\mathcal{N}(\mu_2,\sigma^2)$ ind\'ependants, avec variance commune $\sigma^2$ inconnue. Un IC pour $\mu_1-\mu_2$ est :
\[
(\xbar-\bar{Y}) \pm t_{\alpha/2,\nu}\cdot S_p\sqrt{\frac{1}{n_1}+\frac{1}{n_2}},
\]
o\`u $S_p^2=\frac{(n_1-1)S_1^2+(n_2-1)S_2^2}{n_1+n_2-2}$ est la variance pool\'ee et $\nu=n_1+n_2-2$.
\end{theoreme}

% ------------------------------------------------------------
\section{Taille d'\'echantillon}
% ------------------------------------------------------------

\begin{proposition}
Pour estimer $\mu$ avec une marge d'erreur $E$ au niveau $1-\alpha$ :
\[
n \geq \left(\frac{z_{\alpha/2}\cdot\sigma}{E}\right)^2.
\]
Pour une proportion : $n\geq \frac{z_{\alpha/2}^2\cdot p(1-p)}{E^2}\leq \frac{z_{\alpha/2}^2}{4E^2}$.
\end{proposition}

% ------------------------------------------------------------
\section{Impl\'ementation Python}
% ------------------------------------------------------------

\begin{python}
\begin{minted}{python}
import numpy as np
from scipy import stats

# --- IC pour la moyenne (variance inconnue) ---
data = np.array([12.1, 11.8, 12.5, 12.3, 11.9, 12.7, 12.0, 12.4, 11.7, 12.2])
n = len(data)
xbar = np.mean(data)
s = np.std(data, ddof=1)
alpha = 0.05
t_crit = stats.t.ppf(1 - alpha/2, df=n-1)

ic_lower = xbar - t_crit * s / np.sqrt(n)
ic_upper = xbar + t_crit * s / np.sqrt(n)
print(f"IC 95% pour mu: [{ic_lower:.4f}, {ic_upper:.4f}]")

# Avec scipy directement
ci = stats.t.interval(1-alpha, df=n-1, loc=xbar, scale=s/np.sqrt(n))
print(f"scipy:          [{ci[0]:.4f}, {ci[1]:.4f}]")

# --- IC pour la variance ---
chi2_lower = stats.chi2.ppf(alpha/2, df=n-1)
chi2_upper = stats.chi2.ppf(1-alpha/2, df=n-1)
ic_var = ((n-1)*s**2/chi2_upper, (n-1)*s**2/chi2_lower)
print(f"IC 95% pour sigma^2: [{ic_var[0]:.4f}, {ic_var[1]:.4f}]")

# --- IC pour une proportion ---
n_sondage, x_fav = 1000, 520
p_hat = x_fav / n_sondage
z_crit = stats.norm.ppf(1 - alpha/2)
margin = z_crit * np.sqrt(p_hat * (1-p_hat) / n_sondage)
print(f"IC 95% pour p (Wald):   [{p_hat-margin:.4f}, {p_hat+margin:.4f}]")

# IC de Wilson
p_tilde = (x_fav + z_crit**2/2) / (n_sondage + z_crit**2)
margin_w = z_crit * np.sqrt(p_tilde*(1-p_tilde)/(n_sondage+z_crit**2))
print(f"IC 95% pour p (Wilson): [{p_tilde-margin_w:.4f}, {p_tilde+margin_w:.4f}]")

# --- Taille d'echantillon ---
E = 0.03  # marge d'erreur souhaitee
n_needed = (z_crit**2 * 0.25) / E**2  # p(1-p) <= 0.25
print(f"Taille minimale pour marge {E} sur proportion: n >= {int(np.ceil(n_needed))}")

# --- Visualisation : couverture de l'IC ---
import matplotlib.pyplot as plt
mu_true = 12.0
sigma_true = 0.3
n_ic = 10
n_rep = 50
np.random.seed(1)
fig, ax = plt.subplots(figsize=(10, 6))
for i in range(n_rep):
    sample = np.random.normal(mu_true, sigma_true, n_ic)
    m, se = np.mean(sample), np.std(sample, ddof=1)/np.sqrt(n_ic)
    t_c = stats.t.ppf(0.975, df=n_ic-1)
    lo, hi = m - t_c*se, m + t_c*se
    color = 'blue' if lo <= mu_true <= hi else 'red'
    ax.plot([lo, hi], [i, i], color=color, lw=1.5)
    ax.plot(m, i, 'o', color=color, markersize=3)
ax.axvline(mu_true, color='green', lw=2, linestyle='--', label=f'$\\mu={mu_true}$')
ax.set_xlabel('Valeur')
ax.set_ylabel('Repetition')
ax.set_title('50 IC a 95% : les rouges ne contiennent pas $\\mu$')
ax.legend()
plt.tight_layout()
plt.savefig("ch05_couverture_ic.pdf")
plt.show()
\end{minted}
\end{python}

% ------------------------------------------------------------
\section{Exercices}
% ------------------------------------------------------------

\begin{exercice}[$\star$ -- IC pratique]
On mesure le poids (en g) de 15 sachets de th\'e : $\bar{x}=2{,}03$, $s=0{,}12$. Construire un IC \`a 99\% pour le poids moyen, en supposant la normalit\'e.
\end{exercice}

\begin{exercice}[$\star$ -- Proportion]
Sur 500 pi\`eces produites, 23 sont d\'efectueuses. Construire un IC \`a 95\% pour le taux de d\'efauts (Wald et Wilson).
\end{exercice}

\begin{exercice}[$\star\star$ -- Taille d'\'echantillon]
On veut estimer la taille moyenne des \'etudiants avec une pr\'ecision de 1~cm au niveau 95\%. On suppose $\sigma\approx 8$~cm. Quelle taille d'\'echantillon faut-il ?
\end{exercice}

\begin{exercice}[$\star\star$ -- IC pour le rapport de variances]
Soient deux \'echantillons ind\'ependants de lois normales. Construire un IC pour $\sigma_1^2/\sigma_2^2$ en utilisant la loi de Fisher.
\end{exercice}

\begin{exercice}[$\star\star\star$ -- Projet : couverture r\'eelle]
Par simulation Monte Carlo, comparer le taux de couverture r\'eel des IC de Wald et de Wilson pour $p=0{,}05$ et $n=20,50,100,500$. Lequel est le plus fiable pour petits $n$ ?
\end{exercice}

% ------------------------------------------------------------
\begin{formules}
\begin{align*}
\text{Moyenne ($\sigma$ connu)} &: \xbar\pm z_{\alpha/2}\frac{\sigma}{\sqrt{n}} \\
\text{Moyenne ($\sigma$ inconnu)} &: \xbar\pm t_{\alpha/2,n-1}\frac{S}{\sqrt{n}} \\
\text{Proportion (Wald)} &: \phat\pm z_{\alpha/2}\sqrt{\frac{\phat(1-\phat)}{n}} \\
\text{Variance} &: \left[\frac{(n-1)S^2}{\chi^2_{\alpha/2}},\;\frac{(n-1)S^2}{\chi^2_{1-\alpha/2}}\right] \\
\text{Taille} &: n\geq\left(\frac{z_{\alpha/2}\sigma}{E}\right)^{\!2}
\end{align*}
\end{formules}
